\chapter{Inverse design for Twistoptics}
\nopagebreak
\label{chap:chapter_5}
\thispagestyle{empty}

%%%%%%%%%%%%%%%%%%%%%%%%%%%%%%%%%%%%%%%%%%%%%%%%%%%%%%%%%%%%%%%%%%%%%%%%
% INTRODUCTION
%%%%%%%%%%%%%%%%%%%%%%%%%%%%%%%%%%%%%%%%%%%%%%%%%%%%%%%%%%%%%%%%%%%%%%%%

\section{Introduction}\label{sec:ch5_intro}

Shrinking optics to the nanometre-scale, overcoming the diffraction limit of light, has been the central goal of nano-optics and nanophotonics for decades. This goal has been boosted by the physics of surface polaritons---light coupled to collective material excitations, such as plasmons or phonons---whose wavelengths can be much shorter than those of light in free-space [REFs]. The ability to control and modulate this light--matter interaction is essential for advancing nano-optics and harnessing its full potential. For instance, when surface polaritons are excited in anisotropic media [REFs], they can exhibit highly directional propagation, making them especially attractive for applications such as sensing [REFs], thermal transport [REFs] and hyperlensing [REFs]. These anisotropic polaritons have been extensively studied in metamaterials [REFs], which, however, present some limitations due to high absorption and fabrication challenges [REFs]. The van der Waals (vdW) materials support anisotropic phonon polaritons (PhPs) [REFs]---light coupled to lattice vibrations---naturally due to their intrinsically asymmetric crystal structure. Moreover, polaritonic vdW materials can be exfoliated in thin layers and then be combined in twisted stacks, enhancing the richness and variety of the resulting PhP propagation. This approach, known as twistoptics [REFs], has given rise to exotic optical phenomena at the nanoscale. For instance, canalized PhP propagation has been observed in twisted bilayers of the biaxial vdW crystal $\alpha\text{-}\mathrm{MoO}_3$, where the energy flow is collimated along a single spatial direction, exhibiting diffractionless propagation with deep-subwavelength confinement [REFs]. By adding more layers [REFs], combining different materials [REFs] or varying the layers thickness [REFs], the possibilities to tune the PhP propagation maximizes. Furthermore, advancements in near-field techniques such as scattering-type scanning near-field optical microscopy (s-SNOM) [REFs] have enabled the real-space visualization of PhPs at the nanoscale, providing direct insights into their propagation properties.

Despite these important findings, achieving the desired propagations of light at the nanoscale remains an unsolved challenge in nano-optics. In the case of twistoptics, although its large number of tunable parameters offers remarkable flexibility and non-unique solutions, at the same time it also increases complexity, making the predictability of polariton propagations difficult. In this sense, the combination of inverse design and neural networks (NNs) offers a robust alternative to manual, trial-and-error processes, enabling the design of intricate systems with far greater efficiency and accuracy [REFs]. This approach has proved very effective at addressing the complexity of very different optical systems, as demonstrated in recent studies on metamaterials [REFs], plasmonic nanostructures [REFs] and thermal emitters [REFs].

In this work, we demonstrate precise control and design over the propagation of polaritons at the nanoscale by combining a NN with twisted vdW multilayers. Specifically, desired propagations of PhPs such as canalization and, hitherto unconsidered, bicanalization and tricanalization (fed into our NN as a simple two-dimensional graph in momentum space) are achieved in predicted twisted $\alpha\text{-}\mathrm{MoO}_3$ bilayer and trilayer stacks in unexplored frequency ranges ($\approx 600\,\mathrm{cm}^{-1}$). In addition, we highlight the practical capability of our approach by showing experimentally how to obtain a desired propagation of PhPs (for example, canalization) in a given $\alpha\text{-}\mathrm{MoO}_3$ multilayer by simply adding an extra $\alpha\text{-}\mathrm{MoO}_3$ layer, whose thickness and relative twist angle are predicted by the NN. This example not only validates the predictive power of our framework but also illustrates its versatility, as it can efficiently address both constrained (with initially fixed parameters) and unconstrained design problems. Moreover, we extend the applicability of our approach by applying it to a variety of polaritonic materials that have been barely explored in twistoptics so far, such as $\mathrm{MoOCl}_2$, $\mathrm{CoTeMoO}_6$ and $\alpha\text{-}\mathrm{V}_2\mathrm{O}_5$ thereby enabling access of exotic polariton propagations within an exceptionally wide spectral range covering visible light, infrared and terahertz (THz) frequencies. Notably, such canalization phenomena arise not only from the excitation of PhPs but also from plasmon polaritons, thereby highlighting the versatility of the NN in adapting to distinct fundamental material properties.

%%%%%%%%%%%%%%%%%%%%%%%%%%%%%%%%%%%%%%%%%%%%%%%%%%%%%%%%%%%%%%%%%%%%%%%%
% INVERSE DESIGN FRAMEWORK
%%%%%%%%%%%%%%%%%%%%%%%%%%%%%%%%%%%%%%%%%%%%%%%%%%%%%%%%%%%%%%%%%%%%%%%%

\section{Inverse design framework}\label{sec:ch5_framework}

 \begin{figure*}
   \centering
   \includegraphics[width=\textwidth]{figures/chapter_5/FIG_1.png}
   \caption{Designing nanoscale light propagation by combining deep learning and twistoptics. a, Schematic representation of the inverse design workflow using a NN in twistoptics. The desired polaritonic propagation in real space is transformed into an IFC (using a fast Fourier transform, FFT) that serves as input1 to the NN (along with an additional input, input2, in which optional parameters can be specified), which predicts the optimal multilayer structure. b, The NN (dashed square) operates as a unified model that can be conceptually divided into two parts: (1) a first part where two types of input (input1 defined by $\vec{k}$ and treated by convolutional layers, blue squares and input2 defined by $\vec{p}$ and treated by dense layers, blue circles) generate a vector $\vec{y}_{0}$ with the length of the parameter space, and (2) a second part consisting of a set of $M$ branches that allow finding multivalued solutions. The overall NN, trained with the customized loss function as described in [REFs], accurately predicts the structure parameters ($\theta$, $t$ and $\omega$) for one or multiple configurations.}
   \label{fig:ch5_FIG_1}
 \end{figure*}

The inverse design approach used in this work is summarized in Fig.~\ref{fig:ch5_FIG_1}, where the main NN goal, identifying the twisted multilayer that generates the desired real-space propagation of surface polaritons, for example PhPs, is schematically shown in the top panel (Fig.~\ref{fig:ch5_FIG_1}(a)). As input to the NN we use an isofrequency curve (IFC) (middle panel), a cross-section of the polariton dispersion at a fixed frequency, since it provides an intuitive, detailed and easily designed description of the PhP propagation in real space. Note that although directly related by a Fourier transform operation, real-space numerical simulations of the PhP propagation are computationally expensive to generate and process, making precise predictions more challenging. Furthermore, using IFCs as input offers the advantage of having approximate analytical expressions, which have been accurately reproduced with numerical results [REFs] and experimental data [REFs].

An illustration of the NN design and performance is shown in Fig.~\ref{fig:ch5_FIG_1}(b), including the inputs to the NN and the final predictions of multiple twisted stacks. The inputs consist of an IFC (input1), treated with one-dimensional convolutional layers (blue squares), together with an optional second input (input2), containing constraint parameters, processed by dense layers (blue circles). The convolutional layers are particularly effective for treating IFCs due to the nature of the data: the input IFC is expressed in polar coordinates, where the wavevector magnitude ($k$) is sampled as a function of the polar angle (Supplementary Section 1), forming a naturally ordered one-dimensional sequence that enables the NN to capture local variations [REFs]. In addition, input2 allows the desired polariton propagation to be restricted to specific values of one or more multilayer structural parameters, such as twist angles, thicknesses or frequency. The remainder of the network is composed of dense layers (vertical columns of grey circles), which converge to an output vector $\vec{y}_{0}$ that captures the parameter space with all relevant physical properties and is referred to as the `latent' layer. In principle, $\vec{y}_{0}$ could generate the final output by providing all necessary multilayer parameters. However, due to the presence of degeneracies, that is multilayer twisted configurations might produce the same desired propagation, the network may encounter ambiguities, which makes the training process difficult. To address this issue, we adopted a NN design that enables multiple outputs, or branches [REFs], by generating $M$ vectors ($\vec{y}_{i}$ for $i=1,2,\ldots,M$) from the latent space, each representing system parameters corresponding to the given IFC. Furthermore, these branches use nonlinear functions [REFs], crucial for capturing the complex relationships within the parameter space considered. Multiple solutions not only allow the NN to perform better but also enhance its richness, offering multiple options when designing a device. As a final step in the prediction process, after the NN training, the optimal parameter set will be chosen by computing the analytical IFC for the parameters proposed by the NN (Supplementary Section 2).

%%%%%%%%%%%%%%%%%%%%%%%%%%%%%%%%%%%%%%%%%%%%%%%%%%%%%%%%%%%%%%%%%%%%%%%%
% PHONON POLARITON PROPAGATION IN TWISTED BILAYERS
%%%%%%%%%%%%%%%%%%%%%%%%%%%%%%%%%%%%%%%%%%%%%%%%%%%%%%%%%%%%%%%%%%%%%%%%

\section{Phonon polariton propagation in twisted bilayers}\label{sec:ch5_bilayers}

 \begin{figure*}
   \centering
   \includegraphics[width=\textwidth]{figures/chapter_5/FIG_2.png}
   \caption{NN predictions of PhP canalization and bicanalization in twisted $\alpha\text{-}\mathrm{MoO}_3$ bilayers. a, Schematic of the $\alpha\text{-}\mathrm{MoO}_3$ bilayer indicating the tunable parameters, [$t_{1}$, $t_{2}$], $\theta_1$. $\theta_1$ is the twist angle between the top and bottom layers taking the [001] crystal direction of the latter as reference system. All angles are defined as positive in the counterclockwise direction. b, Learning curve displaying the validation loss. Arrow indicates removal of the less successful branches as epochs increase. c, Predictions of the NN for 2 desired canalization propagations (parallel blue curves) at a selected frequency of $660\,\mathrm{cm}^{-1}$. The resulting analytical IFCs (green lines) are obtained using the parameters supplied by the NN (green numbers). d, IFCs obtained numerically (TMM, displaying the imaginary part of Fresnel reflection coefficient, $\operatorname{Im}(r_p)$) and their real-space simulations in the insets using the parameters supplied by the NN, corroborating the analytical results (overlaid green curves). e, Desired IFCs (blue curves) with non-physical solution in bilayer $\alpha\text{-}\mathrm{MoO}_3$ homostructures and their real-space illustrations as insets (unmatched predictions are in red).}
   \label{fig:ch5_FIG_2}
 \end{figure*}

To test our NN, we will initially focus on $\alpha\text{-}\mathrm{MoO}_3$. The strongly anisotropic crystal structure of this vdW material is reflected in its optical properties, defined by a diagonal permittivity tensor that exhibits very different values of its three components: $\varepsilon_x$, $\varepsilon_y$ and $\varepsilon_z$ [REFs]. In fact, in the far-infrared regime, $\alpha\text{-}\mathrm{MoO}_3$ exhibits up to three spectral bands (the so-called Reststrahlen bands (RBs) [REFs]), in which in addition to different values, the sign of at least one of these permittivity components is negative, thus allowing the excitation of surface PhPs [REFs]. Notably, two of these RBs support in-plane hyperbolic PhPs (RB1 between $544\,\mathrm{cm}^{-1}$ and $850\,\mathrm{cm}^{-1}$ and RB2 between $821\,\mathrm{cm}^{-1}$ and $963\,\mathrm{cm}^{-1}$), enabling twistoptics using twisted $\alpha\text{-}\mathrm{MoO}_3$ multilayers. By adjusting the twist angles and layer thicknesses in these multilayers, transitions between in-plane hyperbolic and elliptical propagation regimes can be observed. More importantly, exotic phenomena, such as canalization [REFs] and ray-like propagation [REFs], emerge at the transitions. While previous research in $\alpha\text{-}\mathrm{MoO}_3$ has primarily focused on the RB2 band, there is growing interest in exploring other spectral regions and potentially expand the frequency range where these exotic optical phenomena occur, such as the RB1, where $\operatorname{Re}(\varepsilon_y)<0$ and $\operatorname{Re}(\varepsilon_x)>0$. This characteristic of the permittivity indicates that RB1 supports in-plane hyperbolic PhPs, which propagate within an angular sector centred along the $y\equiv[001]$ in-plane crystal direction. Accordingly, PhPs in RB1 are candidates to explore twistoptics at an unexplored frequency range and, more interestingly, along complementary in-plane directions compared to those study in RB2. By leveraging our NN model, we aim to predict and analyse the PhP propagation in twisted multilayers within the RB1, thus demonstrating the NN predictive capability in a new physical context.

We begin our analysis with the simplest multilayer: a stack consisting of two thin $\alpha\text{-}\mathrm{MoO}_3$ layers twisted relative to each other by an angle $\theta$ and respective thicknesses $t_{1}$ and $t_{2}$, as schematically represented in Fig.~\ref{fig:ch5_FIG_2}(a). In this configuration, the coordinate axes are aligned with the principal crystal axes of the bottom layer; specifically, the $x(y,z)$ axes correspond to the [100] ([001], [010]) directions of the bottom $\alpha\text{-}\mathrm{MoO}_3$ layer. On the other hand, the twist angle $\theta$ refers to the rotation of the top layer relative to the [001] direction of the bottom $\alpha\text{-}\mathrm{MoO}_3$ layer.

We train the NN using a bilayer dataset calculated analytically, choosing the frequency of the IFC as a selected optional parameter ($\vec{x}=\vec{k},\omega$) (Supplementary Section 2). During training, the number of epochs (complete passes through the dataset) and hyperparameters (preset configurations of the learning algorithm) were carefully monitored by evaluating a loss function designed for the use of branches [REFs] (Fig.~\ref{fig:ch5_FIG_2}(b)). The NN is initially trained with a fixed number of ten branches. During this phase, the activity of each branch is monitored, where `activity' is defined as the number of times the branch provides the correct solution. Once the model converges, the least active branch is removed, and the model is retrained with one fewer branch. The resulting NN learning curve is shown in Fig.~\ref{fig:ch5_FIG_2}(b), showing a rapid decrease with the number of epochs and sharp peaks resulting from the removal of branches. Even when only one of the branches remains ($M=1$), the loss recovers its baseline value, indicating that the NN accuracy is preserved and that degeneracy plays a minor role in bilayers.

% SOURCE NOTE: The PDF lists t11, t21, theta1 and t21, t22, theta2. The repeated t21 is retained; verify layer indexing before finalizing.
Once trained, we use our NN to predict the existence of canalization of PhPs, a paradigmatic phenomenon in twistoptics. Our inverse design approach streamlines this process by defining two parallel lines as the target IFC. This geometry defines canalization, since the normal to the lines (the energy flux direction) is the same for all possible wave vectors. Notably, the desired canalization direction and the minimum PhP confinement can be readily modified by rotating the input lines and adjusting their separation, respectively. Specifically, we conducted our study by considering a fixed line separation of 4.5 $k_0$ and rotation angles from $0^\circ$ to $180^\circ$, covering all possible canalization directions. The NN gives two bilayer configurations with parameters $t_{11}$, $t_{21}$, $\theta_1$ and $t_{21}$, $t_{22}$, $\theta_2$, which lead to perfect canalization. This is corroborated in Fig.~\ref{fig:ch5_FIG_2}(c) by obtaining a perfect match between the analytically obtained IFCs for the resulting NN parameters (green lines) and the target IFCs (blue lines). To further corroborate this result, we also calculate numerically (transfer-matrix method, TMM [REFs]) the IFCs for the predicted parameters, observing again a good match with the target IFCs (Fig.~\ref{fig:ch5_FIG_2}(c), right panels). Its numerically simulated real-space propagations are included in the inset of the TMM figures, providing a more intuitive visualization. These results confirm that the NN accurately predicts the bilayer parameters necessary for achieving a physically feasible IFC target. However, there are directions of canalization that are forbidden, for example, along the $y$ direction (top panel in Fig.~\ref{fig:ch5_FIG_2}(e)). Also, by testing the NN for other exotic propagations, such as canalization along two specific directions, that is, bicanalization (bottom panel in Fig.~\ref{fig:ch5_FIG_2}(e)), our NN does not provide suitable bilayer configurations. These unattainable propagations highlight the limitations of bilayer structures in twistoptics (although this study focuses on RB1, previous studies focused on RB2 suggest that canalization in twisted bilayers is indeed limited to certain directions [REFs]).

%%%%%%%%%%%%%%%%%%%%%%%%%%%%%%%%%%%%%%%%%%%%%%%%%%%%%%%%%%%%%%%%%%%%%%%%
% PHONON POLARITON PROPAGATION IN TWISTED TRILAYERS
%%%%%%%%%%%%%%%%%%%%%%%%%%%%%%%%%%%%%%%%%%%%%%%%%%%%%%%%%%%%%%%%%%%%%%%%

\section{Phonon polariton propagation in twisted trilayers}\label{sec:ch5_trilayers}

 \begin{figure*}
   \centering
   \includegraphics[width=0.95\textwidth]{figures/chapter_5/FIG_3.png}
   \caption{NN predictions of PhP canalization, bicanalization and tricanalization in twisted $\alpha\text{-}\mathrm{MoO}_3$ trilayers. a, Schematic of the trilayer homostructure indicating the tunable parameters ([$t_{1}$, $t_{2}$, $t_{3}$], [$\theta_1$, $\theta_2$]), setting the frequency as a fixed parameter input2. $\theta_1$ and $\theta_2$ are the twist angles of the middle and top layers with respect to the [001] crystal direction of the bottom layer, respectively. All angles are defined as positive in the counterclockwise direction. The thicknesses of the top, middle and bottom layers are denoted by $t_{1}$, $t_{2}$ and $t_{3}$, respectively. b, Learning curve displaying the validation loss value. The arrow indicates removal of less successful branches as epochs increase, observing an increase in the loss function as the number of branches decreases. c, Dependence of the loss jump ($\Delta L$) on the number of branches. A rapid increase can be observed starting from the removal of branch 6, confirming the need of multiple outputs. d, Left, Predictions (pred.) of the NN for two desired canalization propagations (parallel blue curves). The resulting analytical IFCs (light blue, orange and green curves) are obtained using three of the best multivalued output parameters (h). These IFCs are corroborated by numerical calculations (colour plots in the middle and right panels showing the two cases with the most different output parameters). e--g, Same as d for propagations that are forbidden when considering bilayers, such as canalization along the $\alpha\text{-}\mathrm{MoO}_3$ [100] crystal direction (e), bicanalization (f) and tricanalization (g). Vectors $\vec{S}_{1}$, $\vec{S}_{2}$ and $\vec{S}_{3}$ in g indicate the three different propagation directions in both the numerical IFC (middle panel) and real-space numerical simulation (right panel). The results in f,g were obtained using a different NN model from a--e with nine branches, leaving the incident frequency unconstrained. h, NN output parameters for the predictions in d,e ([$t_{1}$, $t_{2}$, $t_{3}$] and [$\theta_1$, $\theta_2$]) and f,g [$t_{1}$, $t_{2}$, $t_{3}$], [$\theta_1$, $\theta_2$] and [$\omega$]).}
   \label{fig:ch5_FIG_3}
 \end{figure*}

One possible solution to expand the number of attainable directions of canalization for PhPs is to consider twisted trilayer homostructures, which have been recently demonstrated to show all-angle canalization tunability [REFs]. This increased versatility can be explained by the more complex coupling between the layers and the presence of more tunable parameters. However, this additional richness considerably complicates the process of predicting desired propagations, making the use of NNs essential for efficiently exploring all possible PhP propagations.

To address this limitation, we trained a NN model with an $\alpha\text{-}\mathrm{MoO}_3$ trilayer database and evaluated its performance with the loss function (learning curve in Fig.~\ref{fig:ch5_FIG_3}(b)) in line with the methodology used for bilayers. The learning curve exhibits an increasing loss value with the removal of branches (in contrast to Fig.~\ref{fig:ch5_FIG_2}(b)), which is clearly observed by the appearance of multiple differentiated steps. This behaviour underscores the need to use multiple branches in trilayer systems due to the presence of degenerate solutions. To highlight the relevance of this outcome, we plot the difference in the loss value for each number of branches (Fig.~\ref{fig:ch5_FIG_3}(c)), observing a rapid increase from the removal of branch 6 and thus indicating the minimum number of branches that should be used to obtain the best results.

Once trained, we first use the NN to predict the canalization of PhPs. As in the case of bilayers, we introduce into the NN a complete set of canalization IFCs as input1 (blue lines). We now obtain multiple solutions that enable canalization (Fig.~\ref{fig:ch5_FIG_3}(d)) not only along the directions accessible with bilayers, but also along previously unachievable ones, such as the $y$ direction (Fig.~\ref{fig:ch5_FIG_3}(e)). In fact, we obtain solutions for all directions (Supplementary Section 4). More importantly, the trilayer model provides multiple trilayer configurations supporting PhP canalization along a specific direction (see Fig.~\ref{fig:ch5_FIG_3}(h) and the left panels of Fig.~\ref{fig:ch5_FIG_3}(d,e) showing the colour-coded analytical IFCs corresponding to the three best predictions). These results are numerically verified, both in momentum and real space (middle and right panels of Fig.~\ref{fig:ch5_FIG_3}(d--f)).

To further explore exotic propagation forms previously unattainable with bilayers, we introduce rhombohedral and hexagonal IFCs into the NN to achieve bicanalization and tricanalization (that is, simultaneous canalization along two or three different in-plane directions), respectively (Fig.~\ref{fig:ch5_FIG_3}(f,g)). In contrast to the bilayer case, we now obtain trilayer configurations that give rise to bicanalization and even to tricanalization (Fig.~\ref{fig:ch5_FIG_3}(f,g)). Moreover, using a NN model with $\vec{x}=(\vec{k})$ as input and ($\vec{t},\vec{\theta},\omega$) as outputs, we observe that bicanalization can be achieved at different frequencies, which clearly shows the great versatility of trilayers. However, for more complex propagation patterns such as tricanalization, the NN identifies only a single viable solution. These results demonstrate that complexity arises not only from the parameter space but also from the richness of the solution space, highlighting the need for optimization-based tools. For completeness, we also evaluated alternative optimization strategies, further demonstrating the robustness and versatility of our NN (Supplementary Section 6).

%%%%%%%%%%%%%%%%%%%%%%%%%%%%%%%%%%%%%%%%%%%%%%%%%%%%%%%%%%%%%%%%%%%%%%%%
% EXPERIMENTAL VALIDATION OF NEURAL-NETWORK-DESIGNED PROPAGATION
%%%%%%%%%%%%%%%%%%%%%%%%%%%%%%%%%%%%%%%%%%%%%%%%%%%%%%%%%%%%%%%%%%%%%%%%

\section{Experimental validation of neural-network-designed propagation}\label{sec:ch5_experiment}

 \begin{figure*}
   \centering
   \includegraphics[width=\textwidth]{figures/chapter_5/FIG_4.png}
   \caption{Real-space nano-imaging of NN-designed PhP propagation in twisted multilayers. a, Top, optical image of the pre-existing twisted bilayer ($t_{1}$ = 291 nm, $t_{2}$ = 208 nm with $\theta_1$ = $130^\circ$) and a schematic of the experimental s-SNOM visualization of PhPs at the surface of a twisted $\alpha\text{-}\mathrm{MoO}_3$ bilayer. The contour of the top and bottom layers is marked in the optical image by black and white dashed rectangles, respectively. Bottom, near-field amplitude of the twisted bilayer, showing hyperbolic propagation of PhPs at an illuminating frequency of $675\,\mathrm{cm}^{-1}$. The inset shows the numerically simulated propagation ($\operatorname{Re}(E_z)$). The white reference axes indicate the optical (crystallographic) directions of the bottom layer. b, Top, schematic of the NN application, fed with the parameters of the bilayer presented in a, for achieving canalization in a twisted trilayer. Bottom, numerically calculated IFCs of the twisted trilayer for the parameters predicted by the NN. Two different desired canalizations (blue lines) are obtained for the following parameters of the third layer: 140 nm in thickness and $96^\circ$ of twist angle (pred. 1) and 280 nm in thickness and $45^\circ$ of twist angle (pred. 2). The IFC of the bilayer is also shown for comparison (yellow lines). c, Same as in a for the predicted trilayer in b (pred. 1). The added third layer is marked in the optical image with a dashed red rectangle. Bottom, near-field amplitude image showing canalization of PhPs along $\varphi_c$ = $158^\circ$ (white dashed lines) in good agreement with the NN prediction and corroborated by numerical simulations (inset).}
   \label{fig:ch5_FIG_4}
 \end{figure*}

Building on the promising results discussed in the previous sections, we aim to experimentally demonstrate the practical use of our NN model and thus highlight its potential applications. In particular, we pursue a desired polaritonic propagation (for example, canalization) in a twisted $\alpha\text{-}\mathrm{MoO}_3$ bilayer (exhibiting hyperbolic propagation) by the simple addition of a third $\alpha\text{-}\mathrm{MoO}_3$ layer.

We start the experimental study by characterizing using s-SNOM (Section~\ref{sec:ch5_methods}) the propagation of PhPs in an $\alpha\text{-}\mathrm{MoO}_3$ bilayer consisting of two superimposed layers (optical image in Fig.~\ref{fig:ch5_FIG_4}(a)). To efficiently launch PhPs, we fabricated an Au nanoantenna on the top layer aligned with its [100] crystal direction. The resulting near-field image (lower panel in Fig.~\ref{fig:ch5_FIG_4}(a)), taken at an incident frequency of $675\,\mathrm{cm}^{-1}$, that is, within RB1, shows concave wavefronts propagating within a hyperbolic sector centred along $-7^\circ$ relative to the [100] crystal axis of the top $\alpha\text{-}\mathrm{MoO}_3$ layer. This PhP propagation is corroborated by performing a numerical simulation (inset). Note, however, that the angle $\varphi$ in the simulations is slightly different from that in the experiment. We attribute it to small inaccuracies in the permittivity used, since RB1 has not yet been explored using near-field techniques [REFs].

After characterizing the hyperbolic PhP propagation in the bilayer at the frequency range of interest, we input the bilayer parameters (input2: thickness, twist angle and frequency) into our NN model to predict an additional third layer (top panel in Fig.~\ref{fig:ch5_FIG_4}(b)) giving rise to canalization. For this aim, we introduce into the NN a complete set of representative IFCs of canalization as input1 (blue lines in Fig.~\ref{fig:ch5_FIG_4}(b)). The NN output gives two trilayer configurations (pred. 1, 140 nm and $96^\circ$; pred. 2, 280 nm and $45^\circ$) that support effective canalization of PhPs. To corroborate these results, we performed numerical simulations (colour plots in Fig.~\ref{fig:ch5_FIG_4}(b)) that show good agreement with the target IFCs (blue lines), especially in the case of pred. 1. Although trilayers can support PhP canalization along arbitrary directions (Fig.~\ref{fig:ch5_FIG_3}), fixing the bilayer configuration and incident frequency restricts the solution to only two directions. On the basis of these results, we fabricated the twisted trilayer corresponding to pred. 1 by adding the predicted $\alpha\text{-}\mathrm{MoO}_3$ layer (optical image in Fig.~\ref{fig:ch5_FIG_4}(c)). The resulting s-SNOM near-field image, taken at $675\,\mathrm{cm}^{-1}$ is shown in Fig.~\ref{fig:ch5_FIG_4}(c) (bottom panel). We observe polaritonic fringes emanating from the Au nanoantenna and propagating constrained to a well-defined direction defined by the in-plane angle $\varphi_c\sim158^\circ$. This result corroborates the emergence of canalization in the trilayer designed by our NN model. Numerical simulations mimicking the experiment show a good agreement (inset). We note, however, a slight difference in the angle of canalization $\varphi_c$, which we ascribe to small discrepancies in the permittivity of $\alpha\text{-}\mathrm{MoO}_3$ in RB1. These results demonstrate the accuracy of the NN in determining optimal system parameters for designing the propagation of PhPs in twisted trilayers. More importantly, we believe that the observation of canalization in the RB1 of $\alpha\text{-}\mathrm{MoO}_3$ is of fundamental relevance, since it opens the door for the study of canalized light--matter interactions at previously unreachable frequencies.

%%%%%%%%%%%%%%%%%%%%%%%%%%%%%%%%%%%%%%%%%%%%%%%%%%%%%%%%%%%%%%%%%%%%%%%%
% EXTENSION TO OTHER POLARITONIC MATERIALS
%%%%%%%%%%%%%%%%%%%%%%%%%%%%%%%%%%%%%%%%%%%%%%%%%%%%%%%%%%%%%%%%%%%%%%%%

\section{Extension to other polaritonic materials}\label{sec:ch5_materials}

 \begin{figure*}
   \centering
   \includegraphics[width=\textwidth]{figures/chapter_5/FIG_5.png}
   \caption{NN predictions of canalized polaritons from visible--near-infrared to terahertz frequencies in twisted multilayers. The top colour bar illustrates the spectral range, from the visible (violet to orange) to the near infrared (NIR) to far infrared (red gradient) and into the terahertz regime (THz) (grey). Below, schematics of the multilayer structures analysed with their NN predicted parameters ([$t_{1}$,$t_{2}$], [$\theta_1$]) that enable canalization at frequencies $\lambda$ = $0.7\,\mu\mathrm{m}$ (visible), $\lambda$ = $10.7\,\mu\mathrm{m}$, $11\,\mu\mathrm{m}$, $15.1\,\mu\mathrm{m}$ and $15.4\,\mu\mathrm{m}$ (far infrared), and $\lambda$ = $32.3\,\mu\mathrm{m}$ (terahertz), respectively. The corresponding analytical IFC (green curve) and the desired canalization propagations (parallel blue curves) are shown in the middle row. Real-space numerically simulated propagation ($\operatorname{Re}(E_z)$) are displayed in the bottom row.}
   \label{fig:ch5_FIG_5}
 \end{figure*}

After validating the predictive capability of our NN for twisted bilayers and trilayers of $\alpha\text{-}\mathrm{MoO}_3$, we study its adaptability to other polaritonic materials. Specifically, we consider twisted homostructures of $\mathrm{MoOCl}_2$ and $\mathrm{CoTeMoO}_6$ ([REFs]), heterostructures composed of $\alpha\text{-}\mathrm{MoO}_3$ and $\alpha\text{-}\mathrm{V}_2\mathrm{O}_5$ ([REFs]), and additional twisted bilayer configurations of $\alpha\text{-}\mathrm{MoO}_3$ (top panel in Fig.~\ref{fig:ch5_FIG_5}). The inclusion of $\mathrm{MoOCl}_2$ allows us to explore the visible and near-infrared spectral ranges, as this material has recently been shown to support plasmon polaritons within these frequencies [REFs]. The remaining materials extend the study into the infrared and terahertz [REFs] frequency domains (by supporting PhPs). As in the previous studies, we focus our analysis on the prediction of canalization phenomena. To this end, the NN was trained using a database of polaritonic IFCs specific to each material configuration (Supplementary Section 5) and subsequently applied using two parallel lines in momentum space as input (blue curves in the middle panel of Fig.~\ref{fig:ch5_FIG_5}). The predicted parameters---including frequency $\omega$, layer thicknesses $t_{1}$ and $t_{2}$, and twist angle $\theta_1$ (depicted in the top panel)---yield analytical IFCs (green lines) that show near-perfect agreement with the target input. These polariton canalization predictions are further corroborated through numerical calculations in real space (bottom panel). Overall, these results establish our NN framework as a versatile and generalizable methodology for the design and prediction of polaritonic propagation, including previously unexplored polariton canalization in terms of frequency range (visible and near infrared), vdW materials ($\mathrm{MoOCl}_2$ and $\mathrm{CoTeMoO}_6$), and material configurations (heterostructures of $\alpha\text{-}\mathrm{MoO}_3$/$\alpha\text{-}\mathrm{V}_2\mathrm{O}_5$).

%%%%%%%%%%%%%%%%%%%%%%%%%%%%%%%%%%%%%%%%%%%%%%%%%%%%%%%%%%%%%%%%%%%%%%%%
% CONCLUSIONS
%%%%%%%%%%%%%%%%%%%%%%%%%%%%%%%%%%%%%%%%%%%%%%%%%%%%%%%%%%%%%%%%%%%%%%%%

\section{Conclusions}\label{sec:ch5_conclusions}

In conclusion, this study demonstrates that combining a NN with multiple twisted polaritonic structures enables desired propagations of light at the nanoscale. Specifically, our NN model accurately predicts twisted $\alpha\text{-}\mathrm{MoO}_3$ bilayers and trilayers supporting exotic PhP propagations, such as canalization, bicanalization and even tricanalization. Furthermore, we showcase its practical applicability by enabling targeted polariton propagation (canalization) in a pre-existing $\alpha\text{-}\mathrm{MoO}_3$ multilayer structure through the simple addition of an extra $\alpha\text{-}\mathrm{MoO}_3$ layer. Beyond solving the inverse design problem---that is, determining specific structural parameters to achieve targeted propagations---the NN can provide multiple solutions in the presence of degeneracies. Notably, the NN can be readily adapted to other materials, enabling us to predict polariton canalization in previously unexplored frequency ranges and material configurations (heterostructures). Overall, our results open the door for the design of multivalued parameter systems exhibiting tailored optical properties at the nanoscale.

%%%%%%%%%%%%%%%%%%%%%%%%%%%%%%%%%%%%%%%%%%%%%%%%%%%%%%%%%%%%%%%%%%%%%%%%
% METHODS
%%%%%%%%%%%%%%%%%%%%%%%%%%%%%%%%%%%%%%%%%%%%%%%%%%%%%%%%%%%%%%%%%%%%%%%%

\section{Methods}\label{sec:ch5_methods}

\subsection{Analytical calculations}

We have used an analytical model [REFs] to calculate the IFC used in the training of the NN for twisted homostructures of biaxial crystals. This analytical equation is a large-$k$-approximation model, considering the finite thickness of the layers and the fields inside them (Supplementary Section 1).

\subsection{Full-wave numerical simulations}

We used the finite boundary element method in full-wave simulations based on commercially available software (COMSOL Multiphysics). A vertically oriented electric dipole placed on top of the upper thin layer acted as a polaritonic launcher. We calculated the vertical component of the electric field ($\operatorname{Re}(E_z)$) at 200 nm on top of the uppermost surface of the twisted $\alpha\text{-}\mathrm{MoO}_3$ bilayers or trilayers. The twisted $\alpha\text{-}\mathrm{MoO}_3$ layers were modelled as biaxial layers rotated with respect to each other on top of a $\mathrm{SiO}_2$ substrate.

\subsection{Transfer-matrix numerical simulations}

The TMM [REFs] was used to obtain the IFCs in Figs.~\ref{fig:ch5_FIG_2}--\ref{fig:ch5_FIG_4}, where the imaginary part of the Fresnel reflection coefficient of the twisted trilayer systems were calculated.

\subsection{Fabrication of twisted trilayer homostructures}

The twisted $\alpha\text{-}\mathrm{MoO}_3$ trilayer homostructure was fabricated using the dry transfer technique [REFs]. First, we performed a mechanical exfoliation from commercial $\alpha\text{-}\mathrm{MoO}_3$ bulk material (Alfa Aesar) using Nitto tape (Nitto Denko Co., SPV 224P). Next, homogeneous $\alpha\text{-}\mathrm{MoO}_3$ flakes were exfoliated from the tape to a transparent poly-(dimethylsiloxane) (PDMS) stamp and selected with the desired thickness by an optical microscope, using a programme that calculates specific thicknesses with remarkable accuracy (the detailed study of this programme is presented in [REFs]). The third $\alpha\text{-}\mathrm{MoO}_3$ flake was precisely aligned with the desired angle respect to the bilayer homostructure on a $\mathrm{SiO}_2$ substrate and release on top of it by heating the PDMS stamp to $100^\circ$.

\subsection{Fabrication of metal antennas}

High-resolution electron beam lithography (100 kV, 100 pA) was performed with submicometre precision on a sample coated with a poly(methyl methacrylate) resist. A standard high-resolution developer (1:3 methyl isobutyl ketone:isopropyl alcohol) was used to pattern the sample, followed by the evaporation of 5 nm Cr and 30 nm Au. A lift-off process defined the antennas. To remove residual organic material, the sample underwent a hot acetone bath ($60\,{}^\circ\mathrm{C}$, 10--15 min), was gently rinsed with isopropyl alcohol (1 min) and dried with nitrogen gas before thermal evaporation. The final gold antenna dimensions were $3.3\,\mu\mathrm{m}$ (length) $\times$ 286 nm (width) $\times$ 40 nm (height).

\subsection{Near-field nano-imaging}

Near-field imaging measurements were performed using a commercial scattering-type scanning near-field optical microscope from Neaspec. The system was equipped with an optical parametric oscillator laser, which covers an emission range from $1.4\,\mu\mathrm{m}$ to $18\,\mu\mathrm{m}$ due to a difference frequency generation stage. Metal-coated (Pt/Ir) atomic force microscopy tips were used, operating at a tapping frequency of $\Omega\approx234\,\mathrm{kHz}$ and with an oscillation amplitude of $\sim$100 nm. The scattered light from the tip was collected using a parabolic mirror and directed into an infrared detector from Kolmar Technologies. To suppress far-field background scattering, the detected signal was demodulated at the third harmonic, while both amplitude and phase signals were extracted using pseudo-heterodyne interferometry.

\section*{Data availability}

All data that support the findings of this study are present in the paper and the Supplementary Information. All raw data generated during the current study are available from the corresponding authors upon request.

\section*{Code availability}

The code used in this paper is available via GitHub at \url{https://github.com/LuciaF-Alvarez-Tomillo/Deep-Learning-and-Twistoptics-Desing.git}.

\section*{Author contributions}

P.A.-G. and L.M.-M. conceived the study. L.F.A.-T., P.C.-B. and S.G.R. designed the NN with input from J.A.-C. and under the supervision of L.M.-M. J.A.-C. and L.F.A.-T. carried out the analytical study and numerical simulations with the help of K.V.V. and A.Y.N. L.F.A.-T. and E.T.-G. performed the sample fabrication and near-field experiments with the help of A.T.M.-L. and P.A.-G. L.F.A.-T., J.A.-C. and P.A.-G. wrote the paper with input from the rest of the authors. P.A.-G. and L.M.-M. supervised the work. All authors contributed to the scientific discussion and paper revisions.

\section*{Funding}

J.A.-C., E.T.-G. and L.F.A.-T. acknowledge support through the Severo Ochoa program from the Government of the Principality of Asturias (grant nos. PA-22-PF-BP21-100, PA-23-PF-BP22-046 and PA-25-PF-BP24-176). P.A.-G. acknowledges support from the European Research Council under Consolidator grant no. 101044461, TWISTOPTICS and the Spanish Ministry of Science and Innovation (State Plan for Scientific and Technical Research and Innovation grant number PID2022-141304NB-I00). L.M.-M., S.G.R. and P.C.-B. acknowledges project nos. PID2023-148359NB-C21 and CEX2023-001286-S (financed by grant no. \url{MICIU/AEI/10.13039/501100011033}) and the Government of Aragon through Project Q-MAD. A.Y.N. acknowledges the Spanish Ministry of Science and Innovation (State Plan for Scientific and Technical Research and Innovation grant number PID2023-147676NB-I00) and the Department of Science, Universities and Innovation of the Basque Government (grant no. PIBA-2023-1-0007 and the IKUR Strategy).

\section*{Competing interests}

The authors declare no competing interests.

\section*{Additional information}

Supplementary information The online version contains supplementary material available at \url{https://doi.org/10.1038/s41563-026-02723-2}.

Correspondence and requests for materials should be addressed to Luis Martín-Moreno or Pablo Alonso-González.

Peer review information Nature Materials thanks Sang-Hyun Oh and the other, anonymous, reviewer(s) for their contribution to the peer review of this work. Peer reviewer reports are available.

Reprints and permissions information is available at \url{https://www.nature.com/reprints}.

\clearpage